\rtmtight
\begin{tblr}{width=\textwidth,colspec={|Q[c,m,2.1cm]|X[l,m]|},hlines,vlines,hspan=minimal,rowsep=1.5pt,colsep=3pt,row{1}={bg=headgray,font=\bfseries}}
\SetCell{c,m}\hfil\logo\hfil & \textbf{POLITEKNIK NEGERI MALANG}\par\textbf{JURUSAN TEKNOLOGI INFORMASI}\par\textbf{PROGRAM STUDI: D4 TEKNIK INFORMATIKA} \\
\SetCell[c=2]{l} \textbf{RENCANA TUGAS MAHASISWA (RTM) DAN RUBRIK PENILAIAN} & \\
Nama Mata Kuliah & Skripsi \\
Kode Mata Kuliah & RTI258001 \quad \textbf{sks / Jam perminggu:} 8 SKS / 16 jam/minggu \quad \textbf{Semester:} 8 \\
Dosen Pengampu & Tim Pengajar Program Studi D4 TI \\
\end{tblr}
\assessmentblock{Proposal Skripsi}{Case Method dan penulisan ilmiah}{SCPMK0802-05803, SCPMK1009-05805}{Mahasiswa menyusun dan mempresentasikan proposal skripsi yang mencakup latar belakang masalah, kajian pustaka, dan metodologi penelitian/pengembangan yang dipilih.}{Mengidentifikasi masalah TIK, mengkaji literatur, merancang metodologi, menyusun naskah proposal, bimbingan dengan pembimbing, dan mempresentasikannya dalam seminar proposal.}{Naskah proposal skripsi lengkap (minimal BAB I--III), slide presentasi seminar proposal, dan berita acara seminar.}{Identifikasi masalah dan relevansi penelitian & 6 & Masalah diidentifikasi secara jelas dari kajian literatur, relevan dengan perkembangan TIK terkini, dan memiliki nilai kontribusi yang terargumentasi. & Masalah teridentifikasi dengan cukup baik, relevansi dengan TIK ada, tetapi kontribusi penelitian masih kurang tajam. & Masalah tidak dirumuskan dengan jelas, relevansi dengan TIK tidak terbukti, atau tidak ada kontribusi yang dapat diidentifikasi. \\ \\
Tinjauan pustaka dan metodologi penelitian & 8 & Tinjauan pustaka komprehensif, sumber primer relevan dan terkini, metodologi dipilih dan dijelaskan dengan tepat sesuai jenis penelitian. & Tinjauan pustaka cukup, tetapi beberapa sumber kurang relevan atau metodologi belum sepenuhnya sesuai. & Tinjauan pustaka sangat terbatas, sumber tidak relevan, atau metodologi dipilih tanpa dasar yang kuat. \\ \\
Kelengkapan dokumen dan presentasi proposal & 6 & Proposal lengkap sesuai format Polinema, disajikan dengan baik dalam seminar, dan dapat menjawab pertanyaan reviewer dengan argumentasi yang kuat. & Proposal cukup lengkap dan presentasi memadai, tetapi beberapa bagian format atau argumentasi saat tanya jawab masih lemah. & Proposal tidak lengkap, tidak sesuai format, atau tidak dapat mempertahankan isi proposal saat tanya jawab. \\}

\assessmentblock{Seminar Kemajuan Penelitian}{Seminar Progress}{SCPMK0403-05801, SCPMK1009-05805}{Mahasiswa mempresentasikan kemajuan penelitian/pengembangan yang telah dicapai, menunjukkan bukti pelaksanaan metodologi, dan mendapatkan masukan untuk perbaikan.}{Melaksanakan penelitian/pengembangan sesuai proposal, mendokumentasikan progress, menyiapkan laporan kemajuan, dan mempresentasikannya dalam seminar.}{Laporan kemajuan penelitian, artefak/data penelitian yang telah dicapai, dan slide presentasi seminar kemajuan.}{Progress pelaksanaan penelitian/pengembangan & 6 & Penelitian/pengembangan telah berjalan sesuai jadwal, target minimal 50\% tercapai, dan bukti pelaksanaan terdokumentasi dengan baik. & Progress ada tetapi belum sesuai jadwal atau target yang direncanakan; sebagian dokumentasi tersedia. & Progress sangat minimal, tidak sesuai rencana, atau tidak dapat menunjukkan bukti pelaksanaan. \\ \\
Kualitas analisis dan temuan sementara & 5 & Temuan sementara dianalisis secara kritis menggunakan kerangka metodologi, insight bermakna dapat diidentifikasi, dan arah penelitian ke depan jelas. & Temuan sementara terdokumentasi tetapi analisisnya masih dangkal atau arah penelitian selanjutnya belum jelas. & Tidak ada temuan sementara yang substansial atau tidak dapat menganalisis data/hasil yang ada. \\ \\
Responsivitas terhadap masukan reviewer & 4 & Masukan reviewer diterima dengan terbuka, pertanyaan dijawab dengan baik, dan rencana tindak lanjut dirumuskan secara konkret. & Sebagian besar masukan ditanggapi dengan baik, tetapi respons terhadap pertanyaan teknis masih terbatas. & Tidak dapat merespons masukan reviewer secara konstruktif atau defensif terhadap kritik. \\}

\assessmentblock{Naskah Skripsi}{Case Method dan penulisan ilmiah}{SCPMK0404-05802, SCPMK0403-05801}{Mahasiswa menyusun naskah skripsi lengkap yang memenuhi standar penulisan ilmiah Polinema, mencakup seluruh proses penelitian/pengembangan dari latar belakang hingga kesimpulan.}{Menyusun seluruh bab skripsi berdasarkan hasil penelitian/pengembangan, bimbingan bertahap dengan pembimbing, revisi berulang, dan pengumpulan naskah final.}{Naskah skripsi lengkap (semua bab) sesuai pedoman Polinema, dilengkapi abstrak Bahasa Indonesia dan Inggris, daftar pustaka, dan lampiran.}{Kelengkapan bab dan standar penulisan ilmiah & 8 & Semua bab tersedia dengan konten lengkap, format sesuai pedoman Polinema, bahasa ilmiah baku, dan sitasi/daftar pustaka menggunakan gaya yang konsisten. & Sebagian besar bab lengkap tetapi format belum sepenuhnya sesuai pedoman, atau ada inkonsistensi penulisan. & Bab penting tidak lengkap, format jauh dari pedoman, atau bahasa tidak memenuhi standar ilmiah. \\ \\
Kualitas hasil penelitian dan analisis & 10 & Hasil penelitian/pengembangan tersaji lengkap, analisis dilakukan secara mendalam, interpretasi didukung teori, dan pembahasan menjawab rumusan masalah. & Hasil penelitian cukup baik tetapi analisis masih terbatas, atau tidak semua rumusan masalah terjawab dengan baik. & Hasil penelitian tidak substansial, analisis sangat dangkal, atau tidak menjawab rumusan masalah. \\ \\
Originalitas dan kontribusi keilmuan & 7 & Penelitian/pengembangan menunjukkan originalitas, memberikan kontribusi nyata terhadap TIK, dan perbedaan dengan penelitian terdahulu dijelaskan dengan jelas. & Ada kontribusi yang dapat diidentifikasi tetapi originalitas atau perbedaan dengan penelitian terdahulu belum cukup eksplisit. & Tidak ada kontribusi yang jelas, atau skripsi hanya mengulang penelitian yang sudah ada tanpa inovasi. \\}

\assessmentblock{Sidang Skripsi}{Presentasi dan defense}{SCPMK0804-05804, SCPMK0404-05802, SCPMK1009-05805}{Mahasiswa mempresentasikan dan mempertahankan skripsi di hadapan tim penguji, menunjukkan penguasaan metodologi, analisis, dan kontribusi penelitian/pengembangan yang dilakukan.}{Menyiapkan presentasi sidang, memaparkan keseluruhan penelitian secara ringkas, menjawab pertanyaan tim penguji, dan merevisi berdasarkan masukan sidang.}{Slide presentasi sidang, naskah skripsi yang telah direvisi pasca sidang, dan berita acara sidang.}{Kualitas presentasi dan penguasaan materi & 12 & Presentasi terstruktur, materi disampaikan dengan jelas dan percaya diri, menunjukkan penguasaan menyeluruh terhadap seluruh isi skripsi. & Presentasi cukup baik tetapi ada bagian yang kurang jelas atau penguasaan beberapa aspek teknis masih terbatas. & Presentasi tidak terstruktur, banyak materi tidak dikuasai, atau penyampaian tidak profesional. \\ \\
Kemampuan mempertahankan argumen ilmiah & 14 & Semua pertanyaan tim penguji dijawab dengan argumentasi ilmiah yang kuat, didukung data/teori, dan menunjukkan pemahaman mendalam terhadap metodologi dan hasil. & Sebagian besar pertanyaan terjawab dengan baik, tetapi beberapa pertanyaan teknis atau metodologis tidak dapat dipertahankan. & Banyak pertanyaan tidak dapat dijawab, argumen lemah, atau tidak dapat mempertahankan validitas hasil penelitian. \\ \\
Validitas solusi dan kontribusi keilmuan & 14 & Solusi yang dihasilkan tervalidasi secara ilmiah, kontribusi terhadap TIK jelas dan signifikan, serta implikasi penelitian untuk pengembangan selanjutnya dirumuskan. & Validasi solusi ada tetapi metode atau cakupannya terbatas; kontribusi keilmuan dapat diidentifikasi meskipun belum optimal. & Solusi tidak tervalidasi atau tidak dapat menjelaskan kontribusi dan relevansi penelitian bagi TIK. \\}

\begin{tblr}{width=\textwidth,colspec={|Q[l,m,3.2cm]|X[l,m]|},hlines,vlines,hspan=minimal,row{1-3}={bg=headgray},rowsep=1.5pt,colsep=3pt}
Jadwal Pelaksanaan: & Minggu 1--17 sesuai RPS. Setiap evaluasi memiliki RTM dan rubrik multi-kriteria. \\
Lain-Lain Yang Diperlukan: & LMS, repositori/naskah, perangkat lunak penelitian, dan perangkat presentasi. \\
Pustaka: & Mengacu pustaka utama dan pendukung pada bagian RPS. \\
\end{tblr}
